Figure \ref{fig:b80-dnp-yield} shows the percentage of the total delayed neutron yield that comes from the six \glspl{DNP} with the largest yield contribution, where the yield from each \gls{DNP} is defined by Equation \eqref{eq:summation-tdny}.
There are 277 \glspl{DNP} in ENDF/B-VIII.0, 59 more than in ENDF/B-VII.1.
Although Figures \ref{fig:b80-dnp-yield} and \ref{fig:b71-dnp-yield} are similar, the main difference is the appearance of $^{86}$As in the top six \gls{DNP} yields.

\begin{figure}
    \centering
    \includegraphics[width=0.5\linewidth]{images/b80/dnp_yield.png}
    \caption{The fraction of the delayed neutron yield from particular \glspl{DNP} from ENDF/B-VIII.0 nuclear data. The six \glspl{DNP} with the largest delayed neutron yield contributions approximately 50\% of the total delayed neutrons per fission event. The remaining 271 \glspl{DNP} provide approximately 50\% of the total delayed neutron yield.}
    \label{fig:b80-dnp-yield}
\end{figure}


Table \ref{tab:b80-yld-hl} shows how the total delayed neutron yield $\nu_d$ and average half-life $\bar{\tau}$ compares between the value from the $I$ \glspl{DNP} and the six \gls{DNP} groups.
The differences between both are well within the uncertainty bounds, indicating the six groups are sufficiently capturing the general behavior of the delayed neutrons.
However, the average half-life differs by approximately 0.3 seconds from Section \ref{sec:b71}, while the total delayed neutron yield differs by approximately 100 $pcm$.
These differences put both results within uncertainty bounds of each other, though the ENDF/B-VIII.0 total delayed neutron yield and average half-life have even larger differences from the values calculated using the IAEA beta-delayed neutron emission database and JENDL-5 than the ENDF/B-VII.1 data.

\begin{table}[]
    \centering
    \caption{ENDF/B-VIII.0 nuclear data delayed neutron yield and average half-life values.}
    \begin{tabular}{lrl}
        \hline
        Parameter & Value & Units\\
        \hline
        $\nu_d(I)$ & $0.0202 \pm 0.0011$ &$ [-]$\\
        $\nu_d(K)$ & $0.0203 \pm 0.0009$ & $[-]$\\
        $\Delta \nu_d$ & $4.64 \pm 142.02$ & $[pcm]$\\
        \hline
        $\bar{\tau} (I)$ & $6.9 \pm 0.3$ & $[s]$\\
        $\bar{\tau} (K)$ & $6.9 \pm 0.3$ & $[s]$\\
        $\Delta \bar{\tau}$ & $0.01 \pm 0.4$ & $[s]$\\
        \hline
    \end{tabular}
    \label{tab:b80-yld-hl}
\end{table}

Table \ref{tab:b80-PCCi} shows the ten largest $PCC_i$ terms, which reveals which \glspl{DNP} have the greatest impact on the \gls{DNP} group parameters $\nu_{d,k}$ and $\tau_k$.
These results are similar to those presented in Section \ref{sec:b71} using ENDF/B-VII.1 data.
The main changes are due to the presence of arsenic-86 and 87, which was not present in the ten largest \glspl{PCC} for ENDF/B-VII.1 data.

\begin{table}[]
    \centering
    \caption{ENDF/B-VIII.0 nuclear data ten largest \glspl{PCC}.}
    \begin{tabular}{ll}
        \hline
       Nuclide & $PCC_{i}$ \\
       \hline
       $^{137}$Sb & 5.052885 \\
       $^{100}$Rb & 4.837795 \\
       $^{137}$I & 4.468315 \\
       $^{87}$Br & 3.350813 \\
       $^{85}$As & 2.685888 \\
       $^{88}$As & 2.107151 \\
       $^{88}$Br & 2.093959 \\
       $^{131}$Cd & 1.988529 \\
       $^{86}$As & 1.862035 \\
       $^{87}$As & 1.472684 \\
        \hline
    \end{tabular}
    \label{tab:b80-PCCi}
\end{table}

Table \ref{tab:b80-Ui} shows the ten largest $U_i$ terms, which corresponds to which \glspl{DNP} have the largest uncertainties relative to their impact on the \gls{DNP} group parameters $\nu_{d,k}$ and $\tau_k$.
The results are very similar to those shown in Table \ref{tab:b71-Ui}, though $^{87}$As has taken the place of $^{102}$Y for having a large \gls{PCC} and relative uncertainty, appearing in both Table \ref{tab:b80-PCCi} and \ref{tab:b80-Ui}.
Another new addition is $^{131}$Cd, which does not have any differences in cumulative fission yield, emission probability, or half-life between ENDF/B-VII.1 and ENDF/B-VIII.0.
Instead the \glspl{PCC} increased because of differences in other \glspl{DNP}, as can be seen by comparing the $^{131}$Cd ENDF/B-VIII.0 $PCC_i$ of 1.989 to the ENDF/B-VII.1 value of 1.889.

\begin{table}[]
    \centering
    \caption{ENDF/B-VIII.0 nuclear data ten largest uncertainty-weighted \glspl{PCC}. Bold \glspl{DNP} are present in Table \ref{tab:b80-PCCi} as well, indicating they have a strong effect on \gls{DNP} group parameters.}
    \begin{tabular}{ll}
        \hline
       Nuclide & $U_{i}$ \\
       \hline
       \textbf{$^{100}$Rb} & \textbf{1.510501} \\
       \textbf{$^{137}$Sb} & \textbf{1.274605} \\
       \textbf{$^{87}$As} & \textbf{0.380737} \\
       $^{99}$Kr & 0.355989 \\
       $^{102}$Y & 0.291115 \\
       $^{84}$As & 0.268318 \\
       \textbf{$^{131}$Cd} & \textbf{0.253599} \\
       $^{98}$Y & 0.245287 \\
       $^{104}$Nb & 0.240678 \\
       $^{124}$Ag & 0.240391 \\
        \hline
    \end{tabular}
    \label{tab:b80-Ui}
\end{table}

Figure \ref{fig:b80-bar-uiv} gives a more detailed view on which \gls{DNP} values specifically lead to the large $U_i$ values.
This offers slightly more insight than Table \ref{tab:b80-Ui}, which sums over the various \gls{DNP} values.
The top four $U_{i,v}$ values are identical as those in Figure \ref{fig:b71-bar-uiv}, indicating nuclear data for $^{100}$Rb and $^{137}$Sb remain problematic in ENDF/B-VIII.0 when calculating \gls{DNP} group parameters.


\begin{figure}
    \centering
    \includegraphics[width=0.5\linewidth]{images/b80/pcc-bar.png}
    \caption{The ten largest uncertainty-weighted \glspl{PCC} $U_{i,v}$ sorted by \gls{DNP} and associated \gls{DNP} value from JENDL-5 cumulative fission yields and IAEA beta-delayed neutron emission database half-lives and emission probabilities.}
    \label{fig:b80-bar-uiv}
\end{figure}


Figure \ref{fig:pcc-charts-b80} shows the $PCC_i$ and $U_i$ plotted on the chart of the nuclides.
These plots do not limit to only the top ten largest values, but instead show how all of the values compare among each other.
The values of $U_i$ are all smaller than those of $PCC_i$, because all of the relative uncertainties are less than 1.
The same points of interest as in Figure \ref{fig:pcc-charts-b71} are still relevant: $^{100}$Rb and $^{137}$Sb.

\begin{figure}[!ht]
\centering
\subfloat[The $PCC_i$, or total \gls{PCC}, calculated for each \gls{DNP}.]{
  \includegraphics[width=0.45\linewidth]{images/b80/chart_PCC.png}
}
\hspace{0.05\linewidth}
\subfloat[The $U_i$, or total uncertainty-weighted \gls{PCC}, for each \gls{DNP}.]{
  \includegraphics[width=0.45\linewidth]{images/b80/chart_PCC_uncertainty.png}
}
\caption{Chart of nuclide plots showing the \glspl{PCC} and uncertainty-weighted \glspl{PCC} using ENDF/B-VIII.0 nuclear data.}
\label{fig:pcc-charts-b80}
\end{figure}

\FloatBarrier